% SPDX-License-Identifier: Apache-2.0
% covers: Syscon
\datasheet{Syscon}{The system control block: identity, build stamp, reset cause, soft reset, and a word that survives a reset.}

\dsfacts{\code{lib/parts/src/syscon.rs}}{\code{txhdl\_parts::syscon::Syscon}}%
{\code{//docs:examples}}

\subsection*{Function}

\code{Syscon} answers three questions only the hardware can answer:
which design is running, why it last restarted, and where to leave a
word for whatever runs next. It also drives the system's reset line.

Four of its words are constants from the type's parameters, so a build
can put a commit and a date into the hardware and a program can read
back which bitstream it is on. Three are registers: the cause, the
reset key and the scratch.

\subsection*{Parameters}

\begin{center}\footnotesize
\begin{tabular}{@{}lp{0.68\textwidth}@{}}
\toprule
Parameter & Meaning \\
\midrule
\code{ID} & The design's identifier, read at \code{0x00}. \\
\code{VERSION} & Its version, read at \code{0x04}. \\
\code{STAMP0} & The build stamp's low word, read at \code{0x08}. A
build would put a commit here. \\
\code{STAMP1} & The build stamp's high word, read at \code{0x0c}. A
build would put a date here. \\
\code{KEY} & The word that must be written to \code{0x14} for a reset
to happen. Any other value is ignored. \\
\bottomrule
\end{tabular}
\end{center}

\subsection*{Register map}

Offsets from the block's base. The block selects the word by bits 4 to
2 of the address and looks at no other bits.

\begin{center}\footnotesize
\begin{tabular}{@{}llp{0.5\textwidth}@{}}
\toprule
Offset & Word & Access \\
\midrule
\code{0x00} & \code{id} & Read only, \code{ID}. \\
\code{0x04} & \code{version} & Read only, \code{VERSION}. \\
\code{0x08} & \code{stamp0} & Read only, \code{STAMP0}. \\
\code{0x0c} & \code{stamp1} & Read only, \code{STAMP1}. \\
\code{0x10} & \code{cause} & Read, and write ones to clear. Bit 0 the
power, 1 the button, 2 the watchdog, 3 software. \\
\code{0x14} & \code{reset} & Write only. Writing \code{KEY} asks for a
reset; anything else does nothing. Reads as zero. \\
\code{0x18} & \code{scratch} & Read and write. A reset does not clear
it. \\
\bottomrule
\end{tabular}
\end{center}

\dstables{Syscon}

\subsection*{Behaviour}

\begin{itemize}
\item Each cycle \code{cause} takes what it held, less any bit a write
to \code{0x10} set to one in that cycle, plus a bit for each reason
asserting now: \code{por}, \code{button}, \code{wdog}, and a write of
\code{KEY} to \code{0x14}.
\item The reasons accumulate. A system reset by the button and then by
software reads both bits until a program clears them, so the case
where two things went wrong is not lost.
\item A write of \code{KEY} to \code{0x14} sets \code{hold} to eight.
While \code{hold} is not zero it counts down by one a cycle. The reset
is held for eight cycles rather than one so that every part of a
design sees it.
\item \code{srst} is high while any of \code{por}, \code{button},
\code{wdog} or the key write is asserting, or while \code{hold} is not
zero. It is combinational from those and from a register.
\item Nothing in this block is cleared by \code{srst}. That is what
lets \code{cause} and \code{scratch} survive a reset, and it is why
the block sits outside whatever \code{srst} drives.
\item A read is answered on \code{bus\_r} in the cycle it is taken, with
\code{OKAY}. A write is taken when its address phase, its beat and
room on \code{bus\_b} are all there.
\end{itemize}

\subsection*{Verification}

\begin{itemize}
\item \code{//lib/examples:ex\_syscon} reads the identity and the
build stamp, reads the power-on cause, leaves a word in
\code{scratch}, clears the cause, asks for a reset with the key, and
checks that the cause comes back as software while the scratch word is
unchanged. It then writes a wrong key and checks that nothing moved,
and presses the button to see the two causes side by side.
\item The build simulates the netlist against that run under nvc and
Verilator: \code{//docs:syscon\_sim\_syscon\_tb\_test} and
\code{//docs:syscon\_vsim\_test}.
\end{itemize}

\subsection*{Used by}

Nothing in the tree yet. Issue 147 asks for it on the board, where the
watchdog of issue 149 would drive \code{wdog} and the bootloader of
issue 135 would use \code{scratch}.

\subsection*{Limits}

The four constants are type parameters, so putting a real commit in
them wants a build that generates the instantiation; this part offers
the words and does not fill them. The reset is eight cycles, not a
parameter. There is no lock, so a program can clear a cause another
program wanted to read. \code{reset} reads as zero rather than as what
was last written.
